\documentclass[10pt,a4paper]{amsart}
\usepackage[margin=19mm]{geometry}
\usepackage{amsmath,amssymb,mathtools}
\usepackage[scr=rsfs,cal=euler]{mathalfa}
\usepackage{xfrac}
\usepackage[unicode,hidelinks,bookmarks=false]{hyperref}
\newcommand{\ie}{i.e.\ }
\newcommand{\cf}{cf.\ }
\newcommand{\resp}{respectively\ }
\newcommand{\Subsection}{\subsection}
\providecommand{\nobreakdash}{\nobreak\hbox{-}}
\setlength{\emergencystretch}{2em}
\multlinegap0pt
\usepackage{scrextend}
\usepackage{epsfig}
\usepackage{forest}
\newtheorem{theorem}{Theorem}
\newtheorem{corollary}[theorem]{Corollary}
\newtheorem{lemma}[theorem]{Lemma}
\newtheorem{proposition}[theorem]{Proposition}
\newtheorem{definition}[theorem]{Definition}
\newtheorem{example}[theorem]{Example}
\newtheorem{conjecture}[theorem]{Conjecture}
\newtheorem{remark}[theorem]{Remark}
\newcommand{\K}{\operatornamewithlimits{\textbf{K}}}
\newcommand{\nobarfrac}{\genfrac{}{}{0pt}{}}
\newcommand{\seqnum}[1]{\href{https://oeis.org/#1}{\rm \underline{#1}}}
\begin{document}
\title[Binary Trees and Sum of Two Squares]{Binary Trees and Sum of Two Squares}
\author{Hongshen Chua}
\date{}
\maketitle
\noindent\emph{Source and licence.} Binary Trees and Sum of Two Squares. Version arXiv:2606.24931v1 (CC BY 4.0). This material is distributed under the Creative Commons Attribution 4.0 International licence, http://creativecommons.org/licenses/by/4.0/, as recorded in the arXiv OAI licence metadata for this exact version and shown on the abstract page https://arxiv.org/abs/2606.24931v1. Full rights notice: licenses/arxiv-2606.24931-CC-BY-4.0.txt.\par\smallskip
\noindent\textit{Editorial scope.} Counted unit: the original \texttt{theorem} environment \texttt{Th: Convergent Matrix} at source lines 162--171 together with its complete proof at lines 173--209; the delivered document keeps source lines 149--210 so that every referenced label and every citation resolves inside it. Native editable source is the paper's own concatenated TeX; the AMS wrapper is editorial formatting only. No publisher class, upstream script or unrelated artwork is redistributed.
\section{Continued Fraction} \label{Ch: Continued Fraction}

Let $ [q_0; q_1, q_2, \ldots] $ represent the continued fraction
\begin{align*}
	q_0 + \cfrac{1}{q_1 + \cfrac{1}{q_2 + \ddots}}.
\end{align*}
The \textit{$ n $-th convergent} $ A_n/B_n $ of a continued fraction is the fraction obtained by truncating the continued fraction at $ q_n $, i.e., $ A_n/B_n = [q_0; q_1, q_2, \ldots, q_n] $. These convergents are related by the following linear recurrences:
\begin{align*}
	A_n = q_n A_{n - 1} + A_{n - 2} \quad \text{ and } \quad B_n = q_n B_{n - 1} + B_{n - 2},
\end{align*}
with the initial values $ A_{-1} = 1 $, $ A_0 = q_0 $, $ B_{-1} = 0 $, and $ B_0 = 1 $.

The matrix binary tree provides a powerful representation of the convergents. Given a continued fraction $ [q_0; q_1, q_2, \ldots] $, we define a path $ W $ through the tree by starting with $ q_0 $ $ R $-moves, followed by $ q_1 $ $ L $-moves, and so on, alternating the direction after each term. By an abuse of notation, we write $ R = \begin{pmatrix} 1 & 1 \\ 0 & 1 \end{pmatrix} $, $ L = \begin{pmatrix} 1 & 0 \\ 1 & 1 \end{pmatrix} $, and the matrix corresponding to the path $ W $ as $ R^{q_0} L^{q_1} \ldots $. This path produces the convergents, as demonstrated in the next theorem.

\begin{theorem} \label{Th: Convergent Matrix} 
Let $ A_n/B_n $ be the $ n $-th convergent of a continued fraction $ [q_0; q_1, q_2, \ldots] $. If a path $ W_n = R^{q_0} L^{q_1} \ldots $ ends with an $ R $-move, then 
\begin{align*}
	W_n = \begin{pmatrix} A_{n - 1} & A_n \\ B_{n - 1} & B_n \end{pmatrix}.
\end{align*}
Otherwise, if $ W_n $ ends with an $ L $-move, then 
\begin{align*}
	W_n = \begin{pmatrix} A_n & A_{n - 1} \\ B_n & B_{n - 1} \end{pmatrix}.
\end{align*}
\end{theorem}

\begin{proof} A quick calculation shows that consecutive $ R $ or $ L $ moves have a simple matrix representation of
\begin{align*}
	R^q = \begin{pmatrix} 1 & q \\ 0 & 1 \end{pmatrix} \quad \text{ and } \quad L^q = \begin{pmatrix} 1 & 0 \\ q & 1 \end{pmatrix}.
\end{align*}

We proceed by induction. For the base case of $ n = 0 $, we have
\begin{align*}
	W_0 = R^{q_0} = \begin{pmatrix} 1 & q_0 \\ 0 & 1 \end{pmatrix} = \begin{pmatrix} A_{-1} & A_0 \\ B_{-1} & B_0 \end{pmatrix}.
\end{align*}
Next, it is simple to verify that
\begin{align*}
	A_1 = q_1 A_0 + A_{-1} = q_0 q_1 + 1 \quad \text{ and } \quad B_1 = q_1 B_0 + B_{-1} = 1.
\end{align*}
Hence, for the case $ n = 1 $, we compute
\begin{align*}
	R^{q_0} L^{q_1} = \begin{pmatrix} 1 & q_0 \\ 0 & 1 \end{pmatrix} \begin{pmatrix} 1 & 0 \\ q_1 & 1 \end{pmatrix} = \begin{pmatrix} q_0 q_1 + 1 & q_0 \\ q_1 & 1 \end{pmatrix} = \begin{pmatrix} A_1 & A_0 \\ B_1 & B_0 \end{pmatrix}.
\end{align*}
This completes the base cases for $ n = 0 $ and $ n = 1 $.

Now, suppose the result holds for all $ n \leq N $ for some $ N \geq 1 $. We prove it for $ n = N + 1 $. If $ W_{N + 1} $ ends with an $ R $-move, then $ W_N $ must have ended with an $ L $-move. By the inductive hypothesis, we get
\begin{align*}
	W_{N + 1} = \begin{pmatrix} A_N & A_{N - 1} \\ B_N & B_{N - 1} \end{pmatrix} \begin{pmatrix} 1 & q_{N + 1} \\ 0 & 1 \end{pmatrix} = \begin{pmatrix} A_N & q_{N + 1} A_N + A_{N - 1} \\ B_N & q_{N + 1} B_N + B_{N - 1} \end{pmatrix}.
\end{align*}
Using the recurrence relation for convergents, the last matrix simplifies to
\begin{align*}
	W_{N + 1} &= \begin{pmatrix} A_N & A_{N + 1} \\ B_N & B_{N + 1} \end{pmatrix}.
\end{align*}
On the other hand, if $ W_{N + 1} $ ends with an $ L $-move, then $ W_N $ must have ended with an $ R $-move, so 
\begin{align*}
	W_{N + 1} = \begin{pmatrix} A_{N - 1} & A_N \\ B_{N - 1} & B_N \end{pmatrix} \begin{pmatrix} 1 & 0 \\ q_{N + 1} & 1 \end{pmatrix} = \begin{pmatrix} q_{N + 1} A_N + A_{N - 1} & A_N \\ q_{N + 1} B_N + B_{N - 1} & B_N \end{pmatrix}.
\end{align*}
Again, the recurrence relation for convergents yields
\begin{align*}
	W_{N + 1} = \begin{pmatrix} A_{N + 1} & A_N \\ B_{N + 1} & B_N \end{pmatrix}.
\end{align*}
This completes the induction step, and hence the proof. 
\end{proof}


\end{document}
